\begin{abstract}
In most message queues one slow job delays every cheap job behind it. Priorities fix this only if someone knows each job's cost and keeps the map current.
\sys{} is a gRPC broker that removes that step: it learns how long each kind of message takes from the acknowledgements it already receives, and serves short expected work first.
On a production-style job queue it cut mean latency by \EOneMeanVsFifoAbs{} against FIFO on AWS and beat a hand-tuned priority map, and as a sidecar the same predictor cut RabbitMQ's mean latency by \EThreeSidecarAbs{}.
Much of this came from plain per-key duration statistics, with no model to train. Under long backlogs, though, the broker's own default aging gave most of the gain back wherever we tested it, and only a design with worker-time weights and pausing aging passed on replayed Azure and Huawei traces, at high load (Azure is in-sample for that design).
Other ideas did not pay off: size-binned keys and learned models added nothing reliable, a queue per job type beat prediction on mean latency, tail latency never improved, and on two HPC logs prediction gave no gain on one and made the other worse.
All pass criteria were frozen before the runs, but the evidence is synthetic or trace-driven: in the trace replays each job simply waits out its recorded duration instead of doing real work, so they test the scheduling but not the contention real handlers would cause. We claim no production readiness.
\end{abstract}
